%=====================================================================
\chapter{Uncertainty and Decision Making}\label{ch:uncertainty}
%=====================================================================
\locator{5}{Part V --- Reasoning Under Uncertainty}

\begin{outcomes}
By the end of this chapter you will be able to:
\begin{tight}
  \item \textbf{Explain} why logic alone cannot represent what a delivery team
        knows, and \textbf{state} probability as a degree of belief.
  \item \textbf{Apply} Bayes' rule to update a belief on evidence, and
        \textbf{avoid} the base-rate fallacy.
  \item \textbf{Use} conditional independence to reduce the size of a joint
        distribution.
  \item \textbf{Choose} an action by maximising expected utility, and
        \textbf{distinguish} a change of belief from a change of decision.
  \item \textbf{Compute} the value of perfect information, and \textbf{say} when
        gathering evidence is worth paying for.
\end{tight}
\end{outcomes}

\begin{prereq}
\chref{propositional} (what a truth value can and cannot express) and
\chref{agents} (the utility-based agent, and rationality as maximising
\emph{expected} performance --- this chapter is where that phrase acquires a
number). School probability; nothing beyond it is assumed.
\end{prereq}

\begin{keyterms}
degree of belief $\bullet$ prior $\bullet$ posterior $\bullet$ likelihood
$\bullet$ base rate $\bullet$ joint distribution $\bullet$ marginalisation
$\bullet$ Bayes' rule $\bullet$ conditional independence $\bullet$ utility
$\bullet$ expected utility $\bullet$ maximum expected utility $\bullet$ decision
table $\bullet$ value of perfect information
\end{keyterms}

\section{What Logic Cannot Say}

\chref{propositional} settled the release gate because everything it was told was
either true or false. Almost nothing a delivery team knows is like that.

\emph{This release probably has a defect.} \emph{The smoke test usually catches
this kind of problem.} \emph{Auth is normally the culprit when payments is slow.}
Every one is actionable, none is a truth value, and writing any of them as a
logical sentence forces a choice between two wrong options: assert it and be
committed to something false, or omit it and reason without what the team knows.

\chref{kr}'s defaults were one answer --- ``normally'', withdrawn on evidence. They
work, and they cannot be combined: two defaults pointing different ways gave
\dsref{diamond}'s conflict, with no principled resolution. Probability is the other
answer, and it is the one that composes.

\begin{definitionbox}[Probability as degree of belief]
$P(a) \in [0,1]$ is an agent's \term{degree of belief} in proposition $a$ given
what it knows --- not a claim about long-run frequency, and not a property of the
world. $P(\mathit{defect}) = 0.1$ says this agent, with this evidence, would accept
either side of a bet at those odds.

A \term{prior} $P(a)$ is the belief before evidence; a \term{posterior}
$P(a \mid b)$ the belief after. The definition of conditional probability,
\[
  P(a \mid b) \;=\; \frac{P(a \wedge b)}{P(b)},
\]
is the whole apparatus, and everything else in this chapter follows from it.
\end{definitionbox}

Unlike defaults, degrees of belief \emph{combine}. Two pieces of evidence pointing
different ways produce a number between them rather than a conflict, and the
combination rule is forced --- not chosen by whoever wrote the system.

\section{Bayes' Rule}

Rearranging the definition gives the tool:
\[
  P(a \mid b) \;=\; \frac{P(b \mid a)\, P(a)}{P(b)}
\]

It matters because the two conditionals are not interchangeable and the available
one is usually the wrong way round. A test report says \emph{how often the test
fires when there is a defect} --- $P(\mathit{alarm} \mid \mathit{defect})$ --- and
what you need is \emph{how likely a defect is, given that it fired} ---
$P(\mathit{defect} \mid \mathit{alarm})$. Confusing the two is the \term{base-rate
fallacy}, and it is the single most consequential probabilistic error made in
engineering organisations.

\begin{examplebox}[The smoke test]
The team's smoke test is a good test:
\begin{tight}
  \item It fires on $90\%$ of releases that really do have a serious defect:
        $P(\mathit{alarm} \mid \mathit{defect}) = 0.9$.
  \item It fires on $20\%$ of releases that do not:
        $P(\mathit{alarm} \mid \neg\mathit{defect}) = 0.2$.
  \item Historically one release in ten has a serious defect:
        $P(\mathit{defect}) = 0.1$.
\end{tight}
The test has just fired. How worried should you be?
\end{examplebox}

\dsfig{%
\begin{tikzpicture}[font=\scriptsize,
  b/.style={rounded corners=3pt, draw=#1, line width=1pt, fill=#1!10,
            text=#1!55!black, font=\scriptsize\bfseries, align=center,
            minimum width=2.5cm, minimum height=0.62cm},
  e/.style={-{Stealth[length=1.8mm]}, neutral!70, line width=0.9pt},
  el/.style={font=\tiny, text=neutral, inner sep=1.5pt}]

\node[b=primary] (all) at (0,3.0) {1000 releases};
\node[b=accent]  (def) at (-3.2,1.6) {100 defective};
\node[b=greenhl] (ok)  at (3.2,1.6)  {900 sound};

\node[b=accent]  (da)  at (-4.8,0.1) {90 alarm};
\node[b=neutral] (dn)  at (-1.7,0.1) {10 silent};
\node[b=goldhl]  (oa)  at (1.7,0.1)  {180 alarm};
\node[b=neutral] (on)  at (4.8,0.1)  {720 silent};

\draw[e] (all) -- node[el, above left] {1 in 10} (def);
\draw[e] (all) -- node[el, above right] {9 in 10} (ok);
\draw[e] (def) -- node[el, left] {90\%} (da);
\draw[e] (def) -- node[el, right] {10\%} (dn);
\draw[e] (ok) -- node[el, left] {20\%} (oa);
\draw[e] (ok) -- node[el, right] {80\%} (on);

\node[rounded corners=3pt, fill=lightbg, draw=primary!35, line width=1pt,
      text=primary!70!black, font=\scriptsize, text width=12.6cm,
      align=flush left, inner sep=6pt] at (0,-1.5)
  {\textbf{Count the alarms.} $90 + 180 = 270$ releases set the test off, and only
   $90$ of them are actually defective:
   \[
     P(\mathit{defect} \mid \mathit{alarm}) \;=\; \frac{90}{270}
     \;=\; \frac{1}{3}
   \]
   \smallskip
   A test that catches $90\%$ of defects has fired, and the release is still
   \emph{twice as likely to be sound as broken}. The $180$ false alarms swamp the
   $90$ true ones, because there are nine times as many sound releases to draw
   them from. That is the base rate doing the work, and it is why the intuitive
   answer --- ``$90\%$'' --- is wrong by a factor of nearly three.};
\end{tikzpicture}}%
{The same calculation as Bayes' rule, in whole releases rather than fractions.
Natural frequencies make the base-rate effect visible, which is why this is the
form to use when explaining a result to a team.}{base-rate}

Silence is far more informative than an alarm here:
$P(\mathit{defect} \mid \neg\mathit{alarm}) = 10/730 = 1/73 \approx 0.014$. A test
with a high false-alarm rate tells you much more by staying quiet.

\section{Conditional Independence}

A \term{joint distribution} over $n$ binary variables has $2^{n}$ entries, which is
\chref{propositional}'s problem again. The escape is the same idea that made
semantic networks tractable: most things are irrelevant to most other things.

\begin{definitionbox}[Conditional independence]
$a$ and $b$ are \term{conditionally independent given} $c$ when
\[
  P(a \wedge b \mid c) \;=\; P(a \mid c)\, P(b \mid c)
\]
Equivalently $P(a \mid b \wedge c) = P(a \mid c)$: once $c$ is known, $b$ says
nothing further about $a$.
\end{definitionbox}

The smoke test and a customer complaint are both caused by the defect, and both are
evidence for it --- but once you \emph{know} whether there is a defect, the
complaint tells you nothing more about whether the test will fire. That is a
conditional independence, and asserting it replaces a table of $2^{3} = 8$ numbers
with three small ones. \chref{bayesnets} makes this the organising principle of a
whole representation.

\section{From Belief to Decision}

A belief is not yet an action. The team does not need to know how likely the defect
is; they need to know whether to ship.

\begin{definitionbox}[Expected utility and the principle of maximum expected utility]
A \term{utility} $U(s)$ is a number expressing how much outcome $s$ is worth. The
\term{expected utility} of action $a$ is
\[
  \mathrm{EU}(a) \;=\; \sum_{s} P(s \mid a)\, U(s)
\]
and a rational agent chooses $a^{*} = \arg\max_{a} \mathrm{EU}(a)$. This is the
precise form of the definition of rationality given informally in \chref{agents}.
\end{definitionbox}

\begin{examplebox}[Ship now or delay a week]
Four outcomes, with the team's own numbers:

\smallskip
\noindent\begin{tabular}{@{}L{3.6cm}L{4.6cm}L{6.0cm}@{}}
\toprule
\textbf{Action} & \textbf{If there is a defect} & \textbf{If there is not} \\
\midrule
Ship now & $-100$ \newline incident, rollback, reputation &
$+40$ \newline a week of revenue earned early \\
Delay a week & $-5$ \newline caught it; a small slip &
$-10$ \newline a week lost for nothing \\
\bottomrule
\end{tabular}

\smallskip
Note that delaying is never \emph{good} --- both its outcomes are negative. It is
insurance, and insurance always costs something.
\end{examplebox}

With $p = P(\mathit{defect})$:
\begin{align*}
  \mathrm{EU}(\text{ship}) &= -100p + 40(1-p) = 40 - 140p \\
  \mathrm{EU}(\text{delay}) &= -5p - 10(1-p) = 5p - 10
\end{align*}
They cross where $40 - 140p = 5p - 10$, that is at
\[
  p^{*} \;=\; \frac{50}{145} \;=\; \frac{10}{29} \;\approx\; 0.345
\]

\dsfig{%
\begin{tikzpicture}
\begin{axis}[
  width=12.4cm, height=6.2cm,
  xlabel={$p = P(\mathrm{defect})$}, ylabel={expected utility},
  xlabel style={font=\scriptsize}, ylabel style={font=\scriptsize},
  tick label style={font=\tiny},
  xmin=0, xmax=1, ymin=-105, ymax=55,
  grid=both, major grid style={gray!18}, minor grid style={gray!8},
  legend style={font=\tiny, at={(0.97,0.97)}, anchor=north east,
                draw=gray!40, fill=white},
  legend cell align=left, domain=0:1, samples=2]

\addplot[primary, line width=1.3pt] {40 - 140*x};
\addlegendentry{$\mathrm{EU}(\text{ship}) = 40 - 140p$}
\addplot[goldhl, line width=1.3pt] {5*x - 10};
\addlegendentry{$\mathrm{EU}(\text{delay}) = 5p - 10$}

\draw[greenhl, line width=1pt, dash pattern=on 3pt off 2pt]
  (axis cs:0.3448,-105) -- (axis cs:0.3448,55);
\node[font=\tiny\bfseries, text=greenhl!50!black, rotate=90, anchor=south]
  at (axis cs:0.3448,-60) {decision flips, $p^{*}=10/29$};

\addplot[only marks, mark=*, mark size=1.8pt, primary]
  coordinates {(0.1,26)} node[above right, font=\tiny, text=primary]
  {prior $0.1$: ship};
\addplot[only marks, mark=*, mark size=1.8pt, accent]
  coordinates {(0.3333,-6.67)} node[above right, font=\tiny, text=accent]
  {after an alarm $1/3$: \emph{still} ship};
\end{axis}
\end{tikzpicture}}%
{The two expected utilities as the probability of a defect varies. The alarm moves
the belief a long way and leaves the decision alone, because $1/3$ falls just short
of the crossing at $10/29$.}{eu-lines}

\section{A Change of Belief Is Not a Change of Decision}

This is the chapter's most useful idea, and \dsref{eu-lines} contains it.

The alarm fires. Belief in a defect more than triples, from $0.1$ to $1/3$. And the
recommendation does not change: at $p = 1/3$, $\mathrm{EU}(\text{ship}) = -6.67$
against $\mathrm{EU}(\text{delay}) = -8.33$, so shipping is still --- barely ---
the better bet. The flip comes at $10/29 \approx 0.345$, and $1/3 \approx 0.333$
falls just short.

Two consequences worth taking seriously.

\textbf{Do not report beliefs to decision makers as though they were decisions.} An
engineer who says ``the probability of a defect has tripled'' has said something
true and nearly useless. The question is which side of $p^{*}$ the number falls on.

\textbf{Attend to the margin.} At $p = 1/3$ the two options differ by $1.67$ utility
units, which is noise. The honest report is not ``ship'' but ``ship, marginally, and
this is now close enough that better evidence would be worth having'' --- which is
exactly what the next section makes precise.

\section{The Value of Information}

If the decision is nearly balanced, finding out is worth a lot. If it is not, it is
worth nothing at all, however interesting the answer.

\begin{definitionbox}[Value of perfect information]
The \term{value of perfect information} about a variable is the increase in
expected utility from learning it before deciding:
\[
  \mathrm{VPI} \;=\; \Bigl[\textstyle\sum_{s} P(s)\, \max_{a} \mathrm{EU}(a \mid s)
  \Bigr] \;-\; \max_{a} \mathrm{EU}(a)
\]
the first term being what you could achieve choosing the best action for each case
once you know which case you are in. It is never negative, and it is an upper bound
on what the evidence is worth paying for.
\end{definitionbox}

For the release decision, $\mathrm{VPI}(p) = 95p$ below the flip point and
$50 - 50p$ above it. Evaluated at the two beliefs of interest:

\smallskip
\noindent\begin{longtable}{@{}L{4.6cm}L{2.6cm}L{2.6cm}L{4.4cm}@{}}
\toprule
\rowcolor{primary!15}
\textbf{Situation} & \textbf{$P(\text{defect})$} & \textbf{Decision} &
\textbf{Value of perfect information} \\
\midrule
\endhead
Before any test & $0.1$ & ship & $9.50$ \\
After the test fires & $1/3$ & ship (marginally) & $\mathbf{31.67}$ \\
After the test is silent & $1/73$ & ship & $1.30$ \\
At the flip point & $10/29$ & indifferent & $32.76$ (the maximum) \\
\bottomrule
\caption{Information is worth most exactly where the decision is hardest, and
nothing at all where it is easy.}\label{tab:vpi}
\end{longtable}

That shape is the point. Before any evidence, certainty would be worth $9.50$.
After an ambiguous alarm --- which changed nothing about what to do --- it is worth
$31.67$, more than three times as much. Information is valuable in proportion to
how close the decision is to flipping, not in proportion to how interesting it is.

\begin{conceptbox}
This is a usable rule for engineering organisations, and it cuts both ways.

\emph{Stop gathering evidence that cannot change the decision.} If the answer is
``ship'' at every plausible value of $p$, another day of investigation has a value
of zero however much everyone would like to know.

\emph{Gather evidence hard when the decision is balanced.} The correct response to
``it is basically a coin flip'' is not to flip the coin but to notice that this is
precisely the circumstance in which finding out pays best.

And notice that both follow from arithmetic rather than from temperament. This is
what it means for an agent to be rational about its own information gathering ---
the third thing \chref{agents} said rationality does not forbid.
\end{conceptbox}

\begin{worked}[Working the alarm through, end to end]
\textbf{Step 1: the posterior.} With $P(d) = 0.1$,
$P(\mathit{al} \mid d) = 0.9$, $P(\mathit{al} \mid \neg d) = 0.2$:
\[
  P(\mathit{al}) = 0.9 \times 0.1 + 0.2 \times 0.9 = 0.09 + 0.18 = 0.27
\]
\[
  P(d \mid \mathit{al}) = \frac{0.9 \times 0.1}{0.27}
  = \frac{0.09}{0.27} = \frac{1}{3}
\]
Check against \dsref{base-rate}: $90$ true alarms out of $270$. The same number,
and the frequency version is the one to put in front of the team.

\smallskip
\textbf{Step 2: the decision.}
\[
  \mathrm{EU}(\text{ship}) = 40 - 140 \times \tfrac{1}{3}
  = 40 - 46.67 = -6.67
\]
\[
  \mathrm{EU}(\text{delay}) = 5 \times \tfrac{1}{3} - 10 = -8.33
\]
Ship, by $1.67$. Both options are now bad --- the alarm has cost the team something
whatever they do --- but shipping is less bad.

\smallskip
\textbf{Step 3: is more evidence worth buying?} With perfect information we would
delay when there is a defect ($-5$ rather than $-100$) and ship when there is not
($+40$ rather than $-10$):
\[
  \text{with information} = \tfrac{1}{3}(-5) + \tfrac{2}{3}(40)
  = -1.67 + 26.67 = 25.00
\]
\[
  \mathrm{VPI} = 25.00 - (-6.67) = 31.67
\]
So certainty is worth $31.67$ units --- against the $10$ units that delaying a week
costs when the release turns out to be sound. \emph{An investigation costing less
than a week's delay is clearly worth doing}, and that conclusion came from the
numbers rather than from anybody's instinct.

\smallskip
\textbf{Step 4: the sensitivity that matters.} Suppose the team disputes the
$-100$. Redo the flip point with shipping a defect costing $-60$:
$\mathrm{EU}(\text{ship}) = 40 - 100p$, which meets $5p - 10$ at
$p^{*} = 50/105 \approx 0.476$. The alarm's $1/3$ is now far from flipping
anything, and the value of information falls accordingly. \emph{The utilities
matter as much as the probabilities}, and they are the part nobody measures ---
which is the honest weakness of this whole apparatus and the reason to write them
down where they can be argued with.
\end{worked}

\begin{pitfall}
\begin{tight}
  \item \textbf{The base-rate fallacy.} $P(\mathit{alarm} \mid \mathit{defect})
        = 0.9$ does not make $P(\mathit{defect} \mid \mathit{alarm})$ anything like
        $0.9$. Here it is $1/3$. Draw the natural-frequency tree.
  \item \textbf{Reporting a belief as though it were a decision.} ``The probability
        tripled'' is true and does not say what to do. Say which side of $p^{*}$ it
        falls.
  \item \textbf{Ignoring the margin.} At $p = 1/3$ the options differ by $1.67$
        units. Report the closeness, not only the winner.
  \item \textbf{Gathering evidence that cannot change the decision.} If the action
        is the same across every plausible probability, the value of information is
        zero.
  \item \textbf{Leaving utilities implicit.} They decide the answer as much as the
        probabilities do, and unlike the probabilities nobody measures them. Write
        them down so they can be disputed.
  \item \textbf{Assuming independence to make the arithmetic easy.} Two symptoms of
        the same cause are \emph{not} independent, and treating them as such
        double-counts the evidence. They may be conditionally independent
        \emph{given} the cause, which is a different claim --- and the one
        \chref{bayesnets} is built on.
\end{tight}
\end{pitfall}

%---------------------------------------------------------------------
\begin{lab}[Go/No-Go Under Uncertainty]
Bayes' rule with exact arithmetic, the decision, the flip point and the value of
information. Every number the chapter quotes is computed here.
\begin{lstlisting}[language=Python]
"""Chapter 14 lab -- ship now or delay a week.

Fractions throughout, so every printed number can be checked by hand
against the natural-frequency tree of Figure 14.1.
"""

from fractions import Fraction as F

# --- the smoke test -----------------------------------------------
P_DEFECT = F(1, 10)          # one release in ten has a serious defect
SENS = F(9, 10)              # P(alarm | defect)
FPR = F(2, 10)               # P(alarm | no defect)

# --- what each outcome is worth -----------------------------------
UTIL = {
    ("ship", "defect"): -100,    # incident, rollback, reputation
    ("ship", "ok"): 40,          # a week of revenue earned early
    ("delay", "defect"): -5,     # caught it; a small slip
    ("delay", "ok"): -10,        # a week lost for nothing
}
ACTIONS = ("ship", "delay")


def posterior(prior, sens, fpr, alarm=True):
    """Bayes' rule. P(defect | alarm) or P(defect | no alarm)."""
    like_d = sens if alarm else 1 - sens
    like_ok = fpr if alarm else 1 - fpr
    return (like_d * prior) / (like_d * prior + like_ok * (1 - prior))


def eu(action, p):
    return UTIL[(action, "defect")] * p + UTIL[(action, "ok")] * (1 - p)


def best(p):
    a = max(ACTIONS, key=lambda x: eu(x, p))
    return a, eu(a, p)


def vpi(p):
    """Value of perfect information about the defect."""
    informed = (p * max(UTIL[(a, "defect")] for a in ACTIONS)
                + (1 - p) * max(UTIL[(a, "ok")] for a in ACTIONS))
    return informed - best(p)[1]


def flip_point():
    """Solve EU(ship) = EU(delay) for p, exactly.

    EU(ship)  = Ud*p + Uo*(1-p)   = Uo + (Ud-Uo)p
    EU(delay) = Vd*p + Vo*(1-p)   = Vo + (Vd-Vo)p
    """
    ud, uo = UTIL[("ship", "defect")], UTIL[("ship", "ok")]
    vd, vo = UTIL[("delay", "defect")], UTIL[("delay", "ok")]
    return F(vo - uo, (ud - uo) - (vd - vo))


if __name__ == "__main__":
    print("=== the smoke test ===")
    print("P(defect)            =", P_DEFECT)
    print("P(alarm | defect)    =", SENS)
    print("P(alarm | no defect) =", FPR)

    post = posterior(P_DEFECT, SENS, FPR, alarm=True)
    quiet = posterior(P_DEFECT, SENS, FPR, alarm=False)
    print()
    print("P(defect | alarm)    =", post, "=", round(float(post), 4))
    print("P(defect | silence)  =", quiet, "=", round(float(quiet), 4))
    assert post == F(1, 3)

    # the natural-frequency check of Figure 14.1
    n = 1000
    defective = n * P_DEFECT                 # 100
    true_alarms = defective * SENS           # 90
    false_alarms = (n - defective) * FPR     # 180
    print()
    print("in %d releases: %d defective, %d true alarms, %d false"
          % (n, defective, true_alarms, false_alarms))
    print("   %d of %d alarms are real = %s"
          % (true_alarms, true_alarms + false_alarms,
             F(int(true_alarms), int(true_alarms + false_alarms))))
    assert F(int(true_alarms), int(true_alarms + false_alarms)) == post
    print("   a 90%-sensitive test fired, and the release is still")
    print("   TWICE as likely sound as broken.")

    # --- the decision ---------------------------------------------
    print()
    print("=== the decision ===")
    print("situation          P(defect)   EU(ship)  EU(delay)  choose")
    print("-" * 60)
    rows = (("no test yet   ", P_DEFECT),
            ("alarm fired   ", post),
            ("test is silent", quiet))
    for label, p in rows:
        a, _ = best(p)
        print("%s %10s %10.2f %10.2f  %s"
              % (label, round(float(p), 4), float(eu("ship", p)),
                 float(eu("delay", p)), a))

    star = flip_point()
    print()
    print("flip point p* =", star, "=", round(float(star), 4))
    assert star == F(10, 29)
    # just either side of it
    assert best(star - F(1, 1000))[0] == "ship"
    assert best(star + F(1, 1000))[0] == "delay"

    # the headline: belief changed a lot, decision did not change
    assert best(P_DEFECT)[0] == "ship"
    assert best(post)[0] == "ship"
    margin = eu("ship", post) - eu("delay", post)
    print()
    print("the alarm moved belief from %s to %s -- more than triple --"
          % (P_DEFECT, post))
    print("and the decision did not change, because %s < %s."
          % (post, star))
    print("the margin is now only %.2f utility units." % float(margin))
    assert float(margin) < 2

    # --- the value of information ---------------------------------
    print()
    print("=== value of perfect information ===")
    for label, p in rows + (("at the flip point", star),):
        print("   %-18s p=%-8s VPI = %6.2f"
              % (label.strip(), round(float(p), 4), float(vpi(p))))
    assert vpi(post) > 3 * vpi(P_DEFECT)
    print()
    print("Information is worth most where the decision is hardest:")
    print("%.2f before the test, %.2f after an ambiguous alarm."
          % (float(vpi(P_DEFECT)), float(vpi(post))))
    print("Delaying a week costs 10 when the release is sound, so an")
    print("investigation cheaper than that is clearly worth doing.")

    # --- and the utilities matter as much as the probabilities ----
    print()
    print("=== if shipping a defect cost 60 rather than 100 ===")
    UTIL[("ship", "defect")] = -60
    star2 = flip_point()
    print("   flip point moves from %s to %s (%.3f)"
          % (star, star2, float(star2)))
    print("   the alarm's 1/3 is now far from flipping anything:",
          best(post)[0])
    print("   VPI after an alarm falls from %.2f to %.2f"
          % (31.67, float(vpi(post))))
    assert star2 > star
    print()
    print("The utilities decide the answer as much as the")
    print("probabilities do -- and they are the part nobody measures.")
\end{lstlisting}
Then change \texttt{FPR} to $0.02$ --- a test that almost never cries wolf --- and
re-run. The posterior after an alarm jumps to about $0.83$, the decision flips to
\emph{delay}, and the value of information collapses. A test is worth having in
proportion to how often its answer changes what you do.
\end{lab}

%---------------------------------------------------------------------
\begin{checkpoint}
\begin{tightnum}
  \item A test that fires on $90\%$ of defective releases has fired, and the
        posterior probability of a defect is $1/3$. Explain to a delivery manager,
        in two sentences and without algebra, why those numbers are consistent.
  \item The alarm moved the belief from $0.1$ to $1/3$ and did not change the
        recommendation. State the general lesson, and say what an engineer should
        report instead of ``the probability tripled''.
  \item The value of perfect information was $9.50$ before the test and $31.67$
        after an ambiguous alarm. Explain why more evidence became \emph{more}
        valuable after evidence arrived.
\end{tightnum}
\end{checkpoint}

%---------------------------------------------------------------------
\begin{chaptersummary}
\begin{tight}
  \item Probability is a \term{degree of belief} given what is known. Unlike
        \chref{kr}'s defaults it \emph{combines}: two pieces of evidence pointing
        different ways give a number between them, by a rule nobody gets to choose.
  \item \term{Bayes' rule} inverts a conditional. The one you have is usually the
        wrong way round, and assuming otherwise is the \term{base-rate fallacy}: a
        $90\%$-sensitive test fired and left $P(\mathit{defect}) = 1/3$.
  \item Natural frequencies --- $90$ true alarms among $270$ --- make the base rate
        visible. Use them when explaining a result to people.
  \item \term{Conditional independence} is what stops a joint distribution growing
        as $2^{n}$, and it is the organising idea of \chref{bayesnets}.
  \item A rational agent maximises \term{expected utility},
        $\mathrm{EU}(a) = \sum_{s} P(s \mid a) U(s)$. This is the precise form of
        \chref{agents}' definition of rationality.
  \item \emph{A change of belief is not a change of decision.} The alarm tripled
        the probability and left the recommendation alone, because $1/3$ falls
        short of the flip point $10/29$. Report which side of $p^{*}$ you are on,
        and how close.
  \item The \term{value of perfect information} is greatest exactly where the
        decision is most finely balanced --- $31.67$ after the ambiguous alarm
        against $9.50$ before it. Gather evidence when it could change the action,
        and stop when it could not.
  \item The utilities decide the outcome as much as the probabilities, and they are
        the part nobody measures. Write them down so they can be argued with.
\end{tight}
\end{chaptersummary}

\begin{reviewq}
  \item \textbf{(MCQ)} The base-rate fallacy is confusing:
        (a)~$P(a)$ with $P(\neg a)$ (b)~$P(a \mid b)$ with $P(b \mid a)$
        (c)~a prior with a utility (d)~independence with conditional independence.
  \item \textbf{(MCQ)} A rational agent chooses the action maximising:
        (a)~probability of success (b)~utility of the best outcome
        (c)~expected utility (d)~the value of information.
  \item \textbf{(MCQ)} The value of perfect information is largest when:
        (a)~one action is clearly best (b)~the decision is finely balanced
        (c)~the prior is zero (d)~the utilities are equal.
  \item \textbf{(MCQ)} $a$ and $b$ are conditionally independent given $c$ when:
        (a)~$P(a \wedge b) = P(a)P(b)$ (b)~$P(a \mid b \wedge c) = P(a \mid c)$
        (c)~$P(a \mid c) = P(a)$ (d)~$P(c \mid a) = P(c \mid b)$.
  \item \textbf{(Short)} Draw the natural-frequency tree for a test with $80\%$
        sensitivity and a $10\%$ false-alarm rate on a $5\%$ base rate, and give
        $P(\mathit{defect} \mid \mathit{alarm})$.
  \item \textbf{(Short)} Explain why the value of information rose after the alarm,
        even though the alarm did not change the recommended action.
  \item \textbf{(Short)} Give one decision from your own experience where more
        investigation had a value of zero, and say how you would have known in
        advance.
  \item \textbf{(Applied)} The team argues that delaying a week should cost $-25$
        when the release turns out sound, not $-10$. Recompute the flip point
        $p^{*}$, say whether the alarm at $p = 1/3$ now changes the decision, and
        compute the new value of perfect information at that point.
\end{reviewq}
